\documentclass[11pt]{article}
\usepackage[margin=1.1in]{geometry}
\usepackage{amsmath,amssymb,amsthm}
\usepackage[T1]{fontenc}
\usepackage{lmodern}
\usepackage{hyperref}
\hypersetup{colorlinks=true,linkcolor=blue,urlcolor=blue,citecolor=blue}
\newtheorem{theorem}{Theorem}
\newtheorem{lemma}[theorem]{Lemma}
\newtheorem{proposition}[theorem]{Proposition}
\theoremstyle{definition}
\newtheorem{definition}[theorem]{Definition}
\setlength{\parskip}{2pt}
\title{A proof of the conjectured linear recurrence for OEIS A244182}
\author{Adrian Perez Fontelles, Independent researcher}
\date{09 September 2026}
\begin{document}
\maketitle
\section{The conjecture}

The Online Encyclopedia of Integer Sequences records \href{https://oeis.org/A244182}{A244182} as

\begin{quote}\itshape Number of length n+4 0..7 arrays with no five consecutive elements with pattern ababa or abbba (possibly a=b) and new values 0..7 introduced in 0..7 order\end{quote}

and carries in its Formula section the line:

\begin{quote}\ttfamily Empirical: a(n) = 16*a(n-1) -57*a(n-2) -146*a(n-3) +471*a(n-4) +2030*a(n-5) +1606*a(n-6) -9926*a(n-7) -40696*a(n-8) -72180*a(n-9) -23242*a(n-10) +262788*a(n-11) +948320*a(n-12) +2028394*a(n-13) +3106669*a(n-14) +3270400*a(n-15) +1237472*a(n-16) -4174278*a(n-17) -13394281*a(n-18) -25599412*a(n-19) -38689072*a(n-20) -49852288*a(n-21) -56526936*a(n-22) -57307380*a(n-23) -52390492*a(n-24) -43372994*a(n-25) -32559403*a(n-26) -22139372*a(n-27) -13590504*a(n-28) -7488422*a(n-29) -3670971*a(n-30) -1580476*a(n-31) -586225*a(n-32) -181844*a(n-33) -44951*a(n-34) -8042*a(n-35) -840*a(n-36)\end{quote}

That line carries no separate signature; the entry itself is by R. H. Hardin (Jun 21 2014).

The word {\itshape Empirical} --- equivalently {\itshape Conjecture} --- is the entry's own: the statement is recorded as fitted to the published terms, and no proof is given there. As of revision 6, last modified Jul 23 2025, the entry records no proof and links to none, so the statement is open. This note settles it.

Written out, the claim is

\begin{multline*} a(n) = 16\,a(n-1) - 57\,a(n-2) - 146\,a(n-3) + 471\,a(n-4) + 2030\,a(n-5) + 1606\,a(n-6) - \\ 9926\,a(n-7) - 40696\,a(n-8) - 72180\,a(n-9) - 23242\,a(n-10) + 262788\,a(n-11) + 948320\,a(n-12) + \\ 2028394\,a(n-13) + 3106669\,a(n-14) + 3270400\,a(n-15) + 1237472\,a(n-16) - 4174278\,a(n-17) - 13394281\,a(n-18) - \\ 25599412\,a(n-19) - 38689072\,a(n-20) - 49852288\,a(n-21) - 56526936\,a(n-22) - 57307380\,a(n-23) - 52390492\,a(n-24) - \\ 43372994\,a(n-25) - 32559403\,a(n-26) - 22139372\,a(n-27) - 13590504\,a(n-28) - 7488422\,a(n-29) - 3670971\,a(n-30) - \\ 1580476\,a(n-31) - 586225\,a(n-32) - 181844\,a(n-33) - 44951\,a(n-34) - 8042\,a(n-35) - 840\,a(n-36). \end{multline*}

stated without a range; it first has meaning at $n = 37$, the least index at which all of $a(n-1),\dots,a(n-36)$ are terms of the sequence.

\section{The object}

The name is read as follows. The reading is not a guess: it is pinned against the entry's own published terms, and against an independent enumeration, before anything is proved (Section 7).

The objects are the words $w_1 w_2 \cdots w_L$ over $\{0,1,\dots,7\}$ of length $L = n + 4$ satisfying the entry's conditions, which are, clause by clause:

\begin{itemize}

\item {\itshape no five consecutive elements with pattern ababa or abbba (possibly a=b)}

\item {\itshape new values 0..7 introduced in 0..7 order}

\end{itemize}

Each of these is a condition on a window of consecutive terms of bounded width, except the labelling clause where present, which says that the first occurrences of the values happen in increasing order.

$a(n)$ is the number of such arrays. The entry's offset is 1, so its term of index $n$ is the count for $n$ rows.

\section{The count is a walk count}

Read the word one letter at a time. Every condition is decided by a window of at most a fixed number of consecutive letters, so it becomes checkable as soon as that many letters have been read; and the labelling clause, where the entry has one, is decided by how many distinct values have been introduced so far, a number between $0$ and 8.

Take as state the last few letters together with that counter. Appending a letter is a transition; a word of length $L$ is a walk of length $L$ from the empty state, and every word arises from exactly one walk. With $M$ the adjacency matrix of the state digraph, $w$ the indicator of the start and $f$ of the accepting states, $a(n) = w^{\mathsf T}M^{\,L}f$, so the count is C-finite.

\section{The degree bound, derived}

The reachable part of that digraph has $8781$ states, $8781$ after trimming. Identifying states with the same number of completions of every length, and then the mirror identification on the edges coming in, leaves $S = 41$ blocks. The count satisfies the linear recurrence given by the characteristic polynomial of the $41\times41$ quotient matrix; that degree bound is computed here, not assumed.

\section{The decision procedure}

Let $c_1,\dots,c_D$ be the coefficients of a claimed recurrence $a(n) = \sum_{i=1}^{D} c_i\,a(n-i)$, let $q(t) = t^{D} - \sum_{i=1}^{D} c_i t^{D-i}$ be its characteristic polynomial, and put

\[u_m \;=\; \tilde w^{\mathsf T} L^{\,m-1} q(L)\,\tilde f \qquad (m \ge 1).\]

Expanding the powers of $L$ and using the displayed formula for $a$, one gets $u_m = a(n) - \sum_{i=1}^{D} c_i\,a(n-i)$ with $n = m + D + 0$. So $u_m$ is exactly the residual of the claim at $n = m + D + 0$, and the recurrence holds at an index $n$ if and only if the corresponding $u_m$ vanishes.

Now $u_m = \tilde w^{\mathsf T} L^{m-1} y$ with $y = q(L)\tilde f$ a fixed vector, so $(u_m)$ satisfies the characteristic polynomial of $L$, of degree $41$: writing that polynomial as $t^{41} - \sum_{i=1}^{41} \lambda_i t^{41-i}$ gives $u_{m+41} = \sum_{i=1}^{41} \lambda_i u_{m+41-i}$. Hence $41$ consecutive vanishing residuals force every later residual to vanish, by induction. The test is therefore finite; and because it locates the {\itshape last} nonvanishing residual, it pins the threshold exactly rather than merely bounding it.

Everything is carried out in exact integer arithmetic on the vector $L^{m-1}y$. The matrix $M$ is never formed.

\section{Results}

\begin{theorem} For A244182, the recurrence of order $36$ displayed in Section 1 holds for every $n \ge 37$, at every index at which it can be stated.\end{theorem}

{\itshape Proof.} The residuals $u_m$ were computed exactly for $m = 1,\dots,102$. Every one of them vanished. After it, more than $S = 41$ consecutive residuals vanish, so by Section 5 every later residual vanishes as well. $\square$

\section{Verification}

Four checks were run, each reproducible from the code accompanying this note, and the first two are what pin the reading of the entry's English.

\begin{itemize}

\item {\bfseries The published terms.} The walk count reproduces all 19 terms the entry publishes, exactly, in integer arithmetic, with the index taken from the entry's stated offset (1) rather than fitted. The first of them are 49, 189, 812, 3832, 19611, 107818, 631917, 3920234,~\dots

\item {\bfseries An independent enumeration.} For $n = 1,\dots,5$ every one of the $8^{1n}$ arrays was written out and tested cell by cell against the condition as the entry states it --- no transfer matrix, no row window, a separate piece of code. The counts agree with the walk count and with the entry.

\item {\bfseries The claim on the raw data.} Each claim was evaluated on the entry's published terms alone, with no model involved, wherever the proved range and the published data overlap. It holds there.

\item {\bfseries The annihilation test.} Carried to the derived bound $S = 41$ over the whole vector, in exact integer arithmetic, never sampled and never truncated.

\end{itemize}

\section{References}

\begin{enumerate}

\item OEIS Foundation Inc., \emph{The On-Line Encyclopedia of Integer Sequences}, entry \href{https://oeis.org/A244182}{A244182}, revision 6, last modified Jul 23 2025.

\item R. P. Stanley, \emph{Enumerative Combinatorics}, Vol.~1, 2nd ed., Cambridge University Press, 2012; \S4.7, the transfer-matrix method.

\item M. Kauers and P. Paule, \emph{The Concrete Tetrahedron}, Springer, 2011; C-finite sequences and their closure properties.

\item J. Berstel and C. Reutenauer, \emph{Noncommutative Rational Series with Applications}, Cambridge University Press, 2011; linear representations of counting automata and their reduction.

\end{enumerate}
\end{document}
